\section{Methodology}\label{sec:method}
% Remaining subsections land in plan Tasks 5-8:
% 3.3 attribution (two mechanisms) ·
% 3.4 leveling-down mechanism · 3.5 trim+lift editor · 3.6 downstream recipe.

\subsection{Problem Formulation}\label{sec:problem}

We study taxi service allocation on a city discretized into $0.01^{\circ}$ grid
cells $g$ and hourly time blocks $t$; a pair $i=(g,t)$ with recorded activity is
an \emph{active unit}. Each active unit carries a demand $D_i$ --- its recorded
pickups --- and a supply $S_i$ --- active-taxi presence, aggregated over the
$5{\times}5$ cell neighborhood to match how a cruising taxi is visible to nearby
demand. Service received relative to demand is the \emph{service ratio}
$Y_i = S_i / \max(D_i, d_0)$, where the demand floor $d_0 = 0.5$ prevents
near-empty cells from producing unbounded ratios.
% src: PAPER/supply-lift/LIFT_ALGORITHM_REFERENCE.md §3 (conventions table)

Each cell additionally carries demographic covariates $x_i$ describing its
neighborhood: housing price, per-capita compensation, and migrant population
share. These covariates resolve at district granularity --- about ten distinct
profiles on the primary city --- which bounds the resolution of any demographic
analysis built on them (\S\ref{sec:objective}).
% src: PAPER/argument/02_datasets.md

The data itself is a corpus $\mathcal{T}$ of real per-driver trajectories, each
a sequence of passenger-seeking states ending in a pickup. The trajectories,
not the aggregate grids, are the editable objects: $D$ and $S$ are
deterministic aggregations of $\mathcal{T}$, so moving where a driver seeks and
picks up moves demand and supply with it.

\textbf{Task.} Given a corpus $\mathcal{T}$ whose service allocation is
demographically inequitable, produce an edited corpus $\mathcal{T}'$ that
(i) is measurably fairer; (ii) remains realistic at the level of individual
trajectories --- a frozen driver-identity discriminator must still recognize
every edited trajectory as its driver; and (iii) transfers its fairness into a
policy trained on it by imitation. FAMAIL is a \emph{data-augmentation} method:
it edits a small, attribution-targeted slice of the real trajectories --- at
most $k = 10{,}000$ of them, each perturbed within an $L_\infty$ ball of radius
$\varepsilon = 2$ grid cells --- and generates nothing.
% src: LIFT_ALGORITHM_REFERENCE.md §3 (EPSILON_BALL, budget)
The edited slice is then upweighted during downstream training so that its
fairness is not averaged away (\S\ref{sec:downstream}).

\subsection{The Fairness Objective}\label{sec:objective}

FAMAIL edits trajectories by gradient ascent on a single differentiable
objective that combines demographic fairness, spatial equity, and realism.

\textbf{Demand-adjusted demographic fairness ($F_{\mathrm{causal}}$, primary).}
Raw parity across neighborhoods is the wrong target: taxi service legitimately
tracks demand, and a business district \emph{should} see more pickups than a
quiet residential block. We therefore measure fairness only in the component of
service that demand does not explain --- the logic of conditional statistical
parity \cite{corbettdavies2017}, which restricts disparities to those not
explained by permitted factors. Stage one fits a flexible power basis $g_0(D)$
from per-unit demand to the service ratio and takes the residual
$R_i = Y_i - g_0(D_i)$; stage two asks how much of that residual the
demographics explain. With $X$ the z-scored demographic design matrix, $H$ the
projection onto $[\mathbf{1}, X]$, and $M$ the centering matrix,
\begin{equation}\label{eq:fcausal}
F_{\mathrm{causal}}
\;=\; \frac{R^{\top}(I-H)\,R}{R^{\top}M\,R}
\;=\; 1 - r^{2}_{\mathrm{demo}} .
\end{equation}
By the Frisch--Waugh--Lovell theorem \cite{frischwaugh1933,lovell1963}, this
residualize-then-project construction recovers exactly the demographic
coefficients of the full multiple regression, so $r^{2}_{\mathrm{demo}}$ is
interpretable as the share of demand-adjusted service inequity attributable to
demographics; measuring fairness by how well a protected attribute predicts an
outcome follows the fairness-as-predictability logic of \cite{feldman2015}.
Boundary cases orient the scale: $R \in \operatorname{span}(X)$ gives
$F_{\mathrm{causal}} = 0$ (fully unfair), and $R \perp X$ gives $1$ (fully
fair).
% src: PAPER/objective-motivation/MOTIVATION.md; PAPER/argument/03_fairness_theory.md

Two caveats are load-bearing. First, $F_{\mathrm{causal}}$ is an
\emph{associational} quantity --- the partial $R^2$ of a cross-sectional
regression on observational district-level demographics, with no
identification and no counterfactual; the ``causal'' in the name is
historical. With about ten district profiles the association is ecological,
and we make no individual-level claims. Second, adjusting for demand treats
recorded demand as exogenous; but demand recorded where service has
historically been thin is itself suppressed by that under-service. We return
to this in \S\ref{sec:leveling}, where the same assumption turns out to bound
what a demand-only editor can achieve.
% src: PAPER/objective-motivation/MOTIVATION.md (associational caveat; demand endogeneity)

\textbf{Spatial equity ($F_{\mathrm{spatial}}$, secondary).} Independent of
demographics, service should also be spread evenly relative to the taxi supply
actually present. We include a Gini-based term over supply-normalized service
rates --- pickups per active taxi and dropoffs per active taxi, per unit ---
the most widely used inequality index in transportation-equity research
\cite{horchergraham2021,karner2024}, in a differentiable pairwise form so it
can enter the gradient objective:
$F_{\mathrm{spatial}} = 1 - \tfrac{1}{2}\bigl(\mathrm{Gini}(\mathrm{DSR}) +
\mathrm{Gini}(\mathrm{ASR})\bigr)$. Being demographic-independent, it acts as
a spatial smoothness regularizer complementing $F_{\mathrm{causal}}$.
% src: PAPER/argument/03_fairness_theory.md

\textbf{Realism guardrail ($F_{\mathrm{fidelity}}$).} An edit that improves
fairness by producing trajectories no real driver would generate is useless
for training a realistic demand model. Because individual mobility signatures
are highly identifying --- the premise of the trajectory-user-linking
literature \cite{gao2017tuler,zhou2018tulvae,miao2020deeptul} --- we score
realism with a frozen driver-identity discriminator of the ST-SiameseNet
family \cite{ren2020stsiamese}: an edited trajectory that still reads as its
driver has preserved human fidelity. The discriminator is trained once and
never updated during editing, avoiding the instability of a live adversarial
game \cite{hoermon2016}. We state its role honestly: under the small, bounded
edits FAMAIL makes, its gradient with respect to the edit is near zero, so it
serves as a realism \emph{guardrail} rather than a driver of edits;
distributional collapse is monitored separately at evaluation time with a
Jensen--Shannon check over aggregate mobility statistics \cite{feng2020simulate}.
% src: PAPER/objective-motivation/MOTIVATION.md (F_fidelity guardrail framing)

\textbf{Scalarization.} The three terms combine linearly,
\begin{equation}\label{eq:objective}
\mathcal{L} \;=\; \alpha_{\mathrm{sp}} F_{\mathrm{spatial}}
+ \alpha_{\mathrm{ca}} F_{\mathrm{causal}}
+ \alpha_{\mathrm{fi}} F_{\mathrm{fidelity}},
\qquad \alpha = (0.2,\; 0.7,\; 0.1).
\end{equation}
% src: PAPER/objective-motivation/MOTIVATION.md ("Why these weights")
The weights were selected under an explicit criterion --- maximize the
demographic-fairness gain subject to the spatial term not degrading
($\Delta F_{\mathrm{spatial}} \ge 0$) --- and the causal-heavy weighting
reflects, rather than forces, where the editable signal lies: the objective's
gradient geometry is dominated by $F_{\mathrm{causal}}$, the spatial gradient
is far smaller and never overrides it, and the fidelity gradient is dormant
under bounded edits. An empirical Pareto sweep over the weight simplex
confirms that the adopted point lies on the
$(\Delta F_{\mathrm{spatial}}, \Delta F_{\mathrm{causal}})$ frontier.
% TODO(alpha-sweep): finalize frontier sentence when 5/5 lands (~2026-07-11 AM);
% partial 4-point finding: frontier flat (dF_causal +0.0217..+0.0226), shipped point on it.
% src: famail_temporal/results/alpha_sweep/summary/ (pending)

\subsection{Attribution: Localizing the Deficit and the Remedy}\label{sec:attribution}

The objective is optimized one trajectory at a time, so FAMAIL must decide
which trajectories are worth a slot in the edit budget. It uses two
attribution mechanisms, one per editing mode (\S\ref{sec:editor}). The first
locates where existing unfairness concentrates --- it attributes the
\emph{deficit}. The second locates where added taxi presence would most
improve fairness --- it attributes the \emph{remedy}.

\textbf{Deficit attribution (drives trim).} Because $M$ and $I-H$ are
idempotent, the explained share $r^{2}_{\mathrm{demo}}$ of
Eq.~\eqref{eq:fcausal} admits an exact per-unit decomposition,
\begin{equation}\label{eq:unit-attr}
r^{2}_{\mathrm{demo}}
\;=\; \sum_i \frac{(MR)_i^{2} - \bigl((I-H)R\bigr)_i^{2}}{R^{\top}MR},
\end{equation}
a mathematically exact partition of the fairness deficit across active units
rather than a heuristic saliency; a signed variant separates over- from
under-served units. The trajectories whose pickups land in the
highest-deficit units are the ones trim edits.
% src: PAPER/argument/03_fairness_theory.md (per-cell attribution)

\textbf{Supply-gradient attribution (drives lift).} The second mechanism asks,
at every active unit simultaneously: \emph{how much fairer would service be
with marginally more taxi presence here?} Instantiating a zero supply
perturbation $\Delta S$ and differentiating the full objective, a single
backward pass answers all $N \approx 34{,}500$ what-ifs at once, yielding a
value-of-presence map over the city. Its $F_{\mathrm{causal}}$ component has
the closed form
\begin{equation}\label{eq:supply-grad}
\frac{\partial F_{\mathrm{causal}}}{\partial S_i}
\;=\; \frac{2}{R^{\top}MR}\;
\frac{\bigl((I-H)R\bigr)_i - F_{\mathrm{causal}}\,(MR)_i}{\max(D_i,\, d_0)},
\end{equation}
against which the automatic gradient is verified.
% src: PAPER/supply-lift/LIFT_ALGORITHM_REFERENCE.md §4.1 (autograd-verified closed form)

To turn the map into candidates, each trajectory's tail --- its last $\ell$
seeking states and pickup --- is rigidly translated by every integer offset
$\delta \in [-\varepsilon, \varepsilon]^2 \setminus \{0\}$, and the linearized
gain of a translation is the presence-mass-weighted change in the
($5{\times}5$ box-summed) gradient over the tail's positions. A trajectory's
score is its best offset, and all ${\approx}95$k trajectories in the corpus
are ranked by score. The screen only \emph{nominates}: every selected
trajectory is subsequently re-optimized individually under the full objective
(\S\ref{sec:editor}), which derives the actual move. In assembling the edit
budget, trim's selection takes precedence and lift fills the remaining slots
with positive-score nominees.
% src: PAPER/supply-lift/LIFT_ALGORITHM_REFERENCE.md §4.2–4.3 (screen; plan assembly)
